\section{Three voters: both transitions, then Fuse}\label{sec:fuse}
This proof ladder starts with the smallest replacement. Let the identities be
$(o,\ell,a,r)$: the retired process, surviving leader, other survivor and fresh
replacement. The complete two-transition schedule is
\[
 (1,1,1,0)\longrightarrow(0,1,1,0)\longrightarrow(0,1,1,1).
\]
The commands that install a configuration are decided under the configuration
governing their slots. For clarity, the calculations below check every row,
including the final row governing subsequent commands; this also covers either
side of each boundary.

\subsection{The base case and the second transition}
For $w:N\to\mathbb N$, write $T(w)=\sum_{n\in N}w(n)$ and
$M_w(R)=\sum_{n\in R}w(n)$. A strict majority satisfies $2M_w(R)>T(w)$.
At the first boundary any old quorum contains two of $o,\ell,a$, while every
intermediate quorum contains both $\ell,a$. Hence they meet. At the second
boundary every intermediate quorum contains $\ell,a$, and any final quorum
contains two of $\ell,a,r$. Hence they meet too. This proves the cross-era
intersection obligation for both transitions, not just the final one.

For the response set $R=\{\ell,a\}$, the masses are $2,2,2$ and the totals
are $3,2,3$. Thus exactly the same two responding identities are a majority
throughout. Adding the replacement's acknowledgement cannot invalidate that
fact. Abstract singleton-intersection witnesses are
$\{o,\ell\}\cap\{\ell,a\}=\{\ell\}$ at the first boundary and
$\{\ell,a\}\cap\{\ell,r\}=\{\ell\}$ at the second.
These geometric witnesses and the responding quorum serve different purposes:
the fused decision certificate uses $\{\ell,a\}$ and does not ask the retired
identity to reply. The existing module \path{ReincarnationThree.lean} records
both intersections and both casting witnesses.

\subsection{What a fused acknowledgement certifies}
Fuse carries a shared header, a count, and the ordered commands for consecutive
slots. FuseOk carries the individual slot acknowledgements; CommitBatch carries
the individual commits. None of these needs a range encoding. Logical slots,
configuration boundaries and their quorum families remain distinct.

Let $F=(b,s,[c_0,\ldots,c_{k-1}])$. An acknowledgement of $F$ certifies that
the sender accepted every specified command at ballot $b$ and its specified
slot. The leader matches the ballot, first slot, count and slot contents to its
pending envelope, counts each identity once, and checks its own eligibility
when adding its own acknowledgement. Let $R$ be that set of acknowledgers.
Then one collection of replies gives acceptance evidence at every packed slot.
It gives decision evidence at slot $i$ precisely when $R\in\QII_i$.

\begin{theorem}[Fuse telescoping]\label{thm:fuse}
Suppose the certified schedule preserves the invariant $R\in\QII_i$ at each
step, the first slot has $R\in\QII_0$, and every member of $R$ acknowledges
all the packed slots. Every packed slot has a phase-two decision certificate.
If the expanded history satisfies Turner's protocol invariants, the ordinary
era-based agreement theorem applies to each certificate.
\end{theorem}
\begin{proof}
Induct on the slot index. The base is the first quorum. Preservation gives
$R\in\QII_{i+1}$ from $R\in\QII_i$. The fused acknowledgements supply every
acceptance required by that quorum at slot $i+1$. Thus every slot has both a
quorum and its acceptance evidence. Expanding the envelope into those logical
events leaves the evidence unchanged; the existing agreement theorem supplies
uniqueness under its history assumptions. In-order application then advances
the committed prefix through the schedule.
\end{proof}
For three voters, the masses above discharge preservation for both steps.
The leader can collect one FuseOk from $a$, complete its own all-slot
acceptance, and issue the two commits in one CommitBatch. This is one successful
request--reply round followed by commit notification, provided the envelope's
acceptance preconditions are already established.

\section{Five voters and arbitrary solver schedules}\label{sec:fuse-five}
\subsection{The first and second transitions}
For the complete six-transition weighted replacement, let $R$ contain any
three surviving original voters and exclude $o$. The initial total is $5$ and
$M(R)=3$. Doubling gives total $10$ and mass $6$: the first boundary preserves
every quorum exactly. Introducing the replacement's unit vote gives total $11$
and mass at least $6$. The strict inequality $12>11$ proves the second
boundary's responding-quorum condition explicitly.

The remaining rows have totals $10,9,10,5$ and response masses at least
$6,6,6,3$. Each inequality stays strict. These numbers cover the full schedule
with doubling and halving, not the shorter five-voter specialisation in
\path{ReincarnationSafety.lean}. Any larger initial responding quorum excluding
the retired identity contains such a triple and therefore also qualifies.
Theorem~\ref{thm:fuse} now telescopes the same replies across all six
transitions. The ordinary cross-era intersections follow from exact scaling
at the ends and the four unit changes between them. There is no need to list
another five-voter replacement table here.

\subsection{The solver's induction invariant}
For arbitrary endpoints, use a response set $R$ that is a strict majority at
both endpoints. Define its margin $\mu_w(R)=2M_w(R)-T(w)$. The constructive
path first decreases coordinates outside $R$, then increases coordinates
inside $R$, moving only towards their targets. These edits increase the
margin. It then decreases coordinates inside $R$ and increases coordinates
outside $R$ towards their targets. Throughout these latter phases the margin
is at least the strictly positive final margin. Thus $\mu_w(R)>0$ everywhere.
Each unit edit also preserves cross-era intersection; exact scaling preserves
quorum families. Zero-weight membership edits have no effect on either proof.

This supplies a useful runtime certificate: compute the path with $R$ as the
availability set, certify the preserved margin, and use the same fused replies
at every slot. The existing operator solver and its general reachability proof
provide the construction; a Fuse integration must retain its response-set
certificate. Merely knowing that some available majority exists in each row
does not establish that every subset which replied first qualifies in all rows.
If only a larger available set $L$ was certified, the leader can use per-slot
quorum accounting on the packed acknowledgements, or wait for a response set
which satisfies the common certificate. Neither requires another wire round
when those responses arrive in the original round.

The necessity is visible in the proposed equal-total swap
$(1,2,1,2)\to(2,1,2,1)$. The response set $\{B,D\}$ has mass $4$ initially
and $2$ finally, with total $6$ at both endpoints. The same replies cannot be
counted as a final majority. Even expanding the swap into legal unit edits
does not alter that endpoint calculation. This is a negative control for the
certificate, not an obstacle to solving the reconfiguration using a suitable
common response set or successive certificates. Phantom identities remain
available for the overlap construction in Appendix~\ref{app:solver}; an
identity which sends no acknowledgement contributes no acceptance evidence.

\section{The atomicity and agreement obligations}\label{sec:fuse-refinement}
The useful state-space reduction is local: one receiver either refuses the
envelope or acknowledges its entire accepted sequence. It removes network
loss between the commands inside that envelope. Different recipients may
still receive different envelopes, and a leader may change while replies or
commits are in flight. The existing agreement proof covers such histories
when its invariants hold.

An indivisible datagram does not itself make execution of the receiver's
handler indivisible. The refinement obligation is therefore to validate the
whole fold before exposing its effects, acknowledge only after all accepts
are recorded, and ensure that interruption before acknowledgement leaves a
state admitted by the ordinary recovery or crash-stop-reincarnation model.
Checksums establish neither a transaction over memory nor a journal
transaction. Under an atomic handler model, expansion is a finite sequence of
ordinary accept actions with no intervening local action. Under a finer model,
its unacknowledged prefixes must also refine ordinary accept histories.

The induction for local acceptance needs its own invariant: the next slot is
consecutive; every command is legal in the accumulated configuration; the
promise still permits its ballot; and the slot's era permits that proposal.
A certified sequence plus preservation of these guards proves that accepting
the first permits accepting the rest. A shared header can hoist repeated
fields out of the wire encoding; it cannot discharge a changing era guard.
In particular, Turner's multi-era argument includes phase-one and value
selection obligations. A Phase2-only envelope must inherit those obligations
from established state or prove the corresponding refinement; the quorum
telescoping theorem alone does not perform phase one.

Concretely, the current ordinary accept path in
\path{src/replica/normal.rs} requires the entry era to equal the ballot era
or its successor. If a shared ballot has era $e$ and expanded entries are
assigned eras $e,e+1,\ldots,e+k$ against that one ballot, this guard implies
$k\leq1$; \path{Fuse.same_ballot_era_guard} checks that arithmetic in Lean.
That is the degenerate reading, not the envelope's. Under the atomic handler
model above, the fold evaluates each payload against the in-memory
configuration the preceding payloads established: payload $i{+}1$'s era
guard is judged in the era payload $i$'s commit folded, exactly as though
the payloads had arrived as separate datagrams with no intervening network.
The entry eras are therefore per-payload facts of the fold, discharged in
order, and the era the header carries disciplines the first payload alone.
Phase-one authority for every packed slot is the header ballot itself:
Lamport's leader executes phase one once for all future instances at the
same proposal number~\cite{fuse-paxos}, which is the induction the
telescoping theorem records.

The proof dependencies are consequently:
\[
\begin{gathered}
\text{legal schedule + preserved response quorum}\\
{}+\text{ all-slot acceptance certificate}\\
\Longrightarrow\text{ decision certificates for all slots};\\
\text{certificates + era history invariants}\\
\Longrightarrow\text{ agreement under reconfiguration}.
\end{gathered}
\]
This is the precise sense in which proving a base case and a preservation rule
replaces a separate proof for every schedule length. It does not assume the
inductive step merely because the first case succeeds.

\subsection{Prior art and executable evidence}
Lamport states the same induction as ``A newly chosen leader executes phase 1 for
infinitely many instances of the consensus algorithm \ldots\ Using the same
proposal number for all instances''~\cite{fuse-paxos}, where the ballot $b$ of a
standard message is defined to apply to all future slots until the leader is
interrupted. The fuse telescoping theorem is the same inductive logic in that the
fuse header ballot number is both a Phase~1 and Phase~2 for all commands in the
fused batch. As the messages are processed as a slab they cannot be interrupted.
Turner supplies the era-history conditions and the weighted reconfiguration
setting~\cite{turner}. The present
telescoping argument combines those obligations with a common response set;
it does not attribute a datagram-atomicity theorem to either source.

The standard-library Python checker
\path{research/weighted-reachability/check_fuse.py} enumerates all identity
subsets, including zero-weight identities. It checks both three-voter
boundaries (64 qualifying cross-quorum pairs), all six five-voter boundaries
(5,760 pairs), and every initial quorum excluding the retired identity
(respectively 2 and 10 sets, counting inclusion of the replacement).
It also checks 29,690 ordered endpoint/response-set combinations for all four
identity profiles in $\{0,1,2\}^4$, covering 99,352 unit edges of the monotone
construction. The equal-total swap is a required negative control. Its output
is \path{research/weighted-reachability/generated/fuse-summary.json}.

These are exhaustive statements about their finite domains. The inductive
argument supplies arbitrary finite schedule length. Lean checks the formal
implication and its stated hypotheses; property tests check the implementation
of packing, validation and commit application. TLA+ can model envelope delivery
as one receiver action after the refinement obligation is discharged.
Maelstrom remains useful for implementation faults involving replies, client
histories, recovery and networking. None of these tools needs to re-prove a
finite arithmetic domain already exhausted by an independent checker, but
each must test the layer for which it supplies evidence.

To replay the arithmetic from the repository root, run \code{python3} on
\path{research/weighted-reachability/check_fuse.py}.
The Lean module is \path{UVRR/Fuse.lean}; its proof
generation and kernel-check evidence are recorded beside the proof ladder.

\section{A phantom slot: phase one of a five-node
reincarnation}\label{sec:phantom}
The telescoping certificate of Section~\ref{sec:fuse} is stated against quorum
families, not against the lifecycle that changes them. This section records the
first phase of a five-node reincarnation in \systemname{}: the initial
configuration with its preallocated phantom slot, the serving leader's declared
overlap, the crash, and the resurrection of the crashed identity into the
phantom slot, with the quorum arithmetic stated at each step. The later phases
--- the forced weight sequence that grants the fresh identity its vote and
retires the old one --- are named where they arise and are not written here.

\begin{definition}[Phantom slot]\label{def:phantom}
A \emph{phantom slot} is a member record of the committed configuration held at
voting weight zero: a preallocated future slot for a resurrection. A phantom
contributes no mass to any quorum; messages from it are discarded by the
membership check, while messages to it are delivered.\footnote{The standby of
the reincarnation design: zero voting weight, streamed the leader's traffic,
and voting only after a committed increment grants it weight.}
\end{definition}

Let $N=\{n_1,\dotsc,n_6\}$ carry the committed weights $(1,1,1,1,1,0)$: five
unit voters $n_1,\dotsc,n_5$ and $n_6$ a phantom slot, preallocated so that a
resurrection finds its record ready, and drawn from the start so that the
figures can name it. The total is $T=5$ and the old view's quorum family is the
strict-majority family of the weights: $R$ qualifies exactly when
$2M_w(R)>5$, that is when $M_w(R)\geq3$ --- a majority of five is three. The
phantom's zero weight decides no membership: $M_w(R\cup\{n_6\})=M_w(R)$ for
every $R$.

The leader is $n_3$, and it declares its overlap as in
Figure~\ref{fig:overlap}: an active decision quorum $q=\{n_1,n_2,n_3\}$ of the
old view and a prospective phase-one quorum $q'=\{n_3,n_4,n_5\}$ for the next,
meeting the casting-vote condition $q\cap q'=\{n_3\}$ of~\eqref{eq:leader}.
While it collects the promises of $q'\setminus\{n_3\}$, the leader services the
current view: its beacon proposal and its own vote are carried forward across
the preparation.\footnote{``Beacon proposal'' is the operator's term for the
leader's continued proposal traffic at the current ballot while its own promise
is withheld; the declared overlap is that proposal together with the leader's
vote.}

Then $n_5$ crashes. A crash changes no weight: only committed configuration
commands do. The committed configuration stays $(1,1,1,1,1,0)$, so the old
view's quorum family is unchanged by the crash; what changes is availability.
Four unit voters survive against one unavailable: the available mass is $A=4$,
the unavailable mass is $D=1$, and $A>D$.

The crashed process restarts dirty and bumps its own identity: it reincarnates
as $n_6=n_5{+}1$, the own-id-increments rule of the reincarnation wire
(\path{Reincarnation.bump} in the Lean development), and announces the pair
$(n_5,n_6)$ to all; only the leader responds. The fresh identity lands in the
preallocated phantom slot --- which is why the slot was preallocated --- and
lights it up as a live standby of weight zero. Until a committed membership
change assigns it weight, $n_6$ votes nowhere.

\begin{proposition}[Phase-one quorum arithmetic]\label{prop:phantom}
Throughout phase one the configured total stays $T=5$ and the majority
threshold stays $3$: the crash removes availability, not weight, and the
phantom adds zero. Explicitly: the old view's quorum family is
$\{R\subseteq N:2M_w(R)>5\}$ computed on $w=(1,1,1,1,1,0)$; the live mass is
$5-1+0=4$, a majority of five minus the crashed vote plus the phantom's zero
weight, and $2\cdot4>5$; any three of the four survivors form a quorum, since
$2\cdot3=6>5$; and the declared overlap pair still meets in exactly
$\{n_3\}$.
\end{proposition}
\begin{proof}
Each clause is arithmetic on the displayed $w$. The family clause restates the
strict-majority definition; a crash edits no weight, so the family is
unaffected. The phantom clause is $w(n_6)=0$. The live-mass clause sums the
four surviving unit weights; the strict inequality $2M>T$ holds at $M=4$ and at
$M=3$. The overlap clause is set intersection:
$\{n_1,n_2,n_3\}\cap\{n_3,n_4,n_5\}=\{n_3\}$. Any other prepare quorum drawn
from the survivors still meets every old decision quorum, because two subsets
each of mass at least $3$ out of a total of $5$ must share a member.
\end{proof}

\begin{figure*}[t]
\centering
\begin{tikzpicture}[x=1cm,y=1cm,>=Stealth,font=\small]
\node at (0,0.45) {$n_1$};
\node at (1.8,0.45) {$n_2$};
\node[fill=green!30,inner sep=2pt] at (3.6,0.45) {leader $n_3$};
\node at (5.4,0.45) {$n_4$};
\node at (7.2,0.45) {$n_5$};
\node[draw,gray,dashed,inner sep=2pt] at (9,0.45) {$n_6$};
\draw[gray] (0,0)--(0,-4.9);
\draw[gray] (1.8,0)--(1.8,-4.9);
\draw[green!55!black] (3.6,0)--(3.6,-4.9);
\draw[gray] (5.4,0)--(5.4,-4.9);
\draw[gray] (7.2,0)--(7.2,-2.2);
\draw[gray,dashed] (7.2,-2.2)--(7.2,-4.9);
\draw[gray,dashed] (9,0)--(9,-3.4);
\draw[gray] (9,-3.4)--(9,-4.9);
\draw[->,green!55!black,dashed] (3.6,-0.3)--(3.6,-1.75);
\draw[dashed,gray] (-0.8,-1.1)--(9.8,-1.1);
\node[anchor=west,align=left,text width=6.9cm] at (10.2,-1.1) {initial
configuration $(1,1,1,1,1,0)$; $n_3$ declares its leader overlap,
$q\cap q'=\{n_3\}$, and services the current view: beacon proposal and own
vote carried forward};
\draw[red,very thick] (7.05,-2.35)--(7.35,-2.05) (7.05,-2.05)--(7.35,-2.35);
\node[anchor=west,align=left,text width=6.9cm] at (10.2,-2.2) {$n_5$ crashes;
a crash changes no weight};
\draw[->,gray,dashed] (7.2,-2.5) .. controls (8.1,-2.95) .. (9,-3.3)
node[midway,above,sloped] {$(n_5,n_6)$};
\draw[dashed,gray] (-0.8,-3.4)--(9.8,-3.4);
\node[anchor=west,align=left,text width=6.9cm] at (10.2,-3.4) {resurrection
announced as $n_6=n_5{+}1$; the preallocated slot lights up at weight $0$};
\node[anchor=west,align=left,text width=6.9cm] at (10.2,-4.55) {phase one
complete: live mass $4$, the old view's family unchanged, $n_6$ a live standby
of weight $0$};
\end{tikzpicture}
\caption{Phase one of a five-node reincarnation with a preallocated phantom
slot. Six membership slots $n_1,\dotsc,n_6$; the phantom $n_6$ is drawn grey
and dashed from the start at weight $0$ and lights up when the crashed $n_5$
resurrects into it. The leader $n_3$ (filled) declares its overlap --- beacon
proposal and own vote carried forward, dashed --- and services the current view
throughout. A crash changes no weight: the configured total stays $5$ and the
majority threshold stays $3$; the live mass is $5-1+0=4$. Time runs downwards
without a scale.}
\label{fig:phantom-phase-one}
\end{figure*}

\begin{remark}[The later phases, named]\label{rem:later-phases}
Phase two is the leader's forced weight sequence over the six-slot membership
of Table~\ref{tab:weights} --- the doubling, the join and increment that grant
$n_6$ its vote, the decrements that retire $n_5$, and the halve back to unit
weights --- with the leader's memo-stream keeping the standby current
throughout. Phase three is the completion: $n_5$ leaves the membership and the
cluster returns to five unit voters, now with $n_6$ among them. Neither is
written here.
\end{remark}
